
\subsubsection{2.1 Reinforcement Learning with Execution Feedback for Code LLMs}

The standard RLEF paradigm assigns binary reward: 1 if all test cases pass, 0 otherwise. CodeRL \citep{Le2022CodeRL} applied this with PPO on the APPS dataset \citep{Hendrycks2021APPS} and reported improvements over supervised fine-tuning on HumanEval \citep{Chen2021HumanEval}. PPOCoder \citep{Shojaee2023PPOCoder} extended the framework with additional reward components (syntax correctness signals) while retaining binary execution reward as the primary signal. RLEF/Gehring et al. [2024] applied binary RLEF to repository-level code repair on SWE-bench. None of these works analyze whether binary reward's group-level gradient properties are appropriate for GRPO, nor do they compare binary reward against ratio reward in controlled ablation.

\subsubsection{2.2 Group Relative Policy Optimization}

GRPO \citep{Shao2024DeepSeekMath} eliminates the critic network by normalizing rewards within a group of completions. This changes the reward sparsity problem qualitatively: where PPO's value function provides a non-zero baseline under sparse execution rewards, GRPO's group normalization redistributes only the signal present within the group --- it cannot create signal where none exists. If all completions in a group receive identical rewards (including all-zero under binary reward with universal failure), the group contributes zero gradient.

DAPO \citep{Yu2025DAPO} is the most prominent GRPO-based code training system. It acknowledges reward sparsity and introduces dynamic sampling temperature and a clip-higher variant. However, DAPO retains binary reward and does not quantify the fraction of training groups receiving zero gradient under early-training conditions. Our analysis shows this fraction is 98.7\% under realistic partial-pass rates.

\subsubsection{2.3 Reward Signal Density and Partial Credit}

Sparse rewards in reinforcement learning have a long history \citep{Sutton2018RL}. Reward shaping \citep{Ng1999Shaping} addresses policy-level sparsity via potential-based bonuses. The ratio reward analyzed here (k/n test cases passing) is a form of partial-credit reward. The contribution of this work is not the reward function itself, but the group-level analysis showing when it provides a mathematical guarantee of gradient signal that binary reward cannot provide, and the precise characterization of the experimental conditions under which this guarantee is active.

\subsubsection{2.4 Positioning}

This work differs from all prior RLEF literature by analyzing gradient signal quality at the GRPO group level --- the atomic unit of computation --- rather than at aggregate training curve level. This framing reveals a failure mode invisible in aggregate metrics and provides a falsifiable condition for when ratio reward outperforms binary reward in terms of gradient signal density.

